% Imporditud: tools/import_eksperimendid.py
% Allikas: Eesti füüsikaolümpiaadi eksperimendi ülesannete
%   kogu (Rael Kalda, 2023), ülesanne Ü38, lahendus L38.
\setAuthor{EFO žürii}
\setRound{piirkonnavoor}
\setYear{1999}
\setNumber{P E1}
\setDifficulty{2}
\setTopic{laineoptika}
\setVahendid{2 ümmargust pliiatsit, kleepriba, küünal või taskulambipirn alusel koos patareiga}
\setPunktid{13}
\setAllikas{Eksperimendi kogu Ü38, lahendus lk 50}
\setLahendusseis{toores}

\prob{Küünlaleek}
Vaadake küünlaleeki või taskulambipirni ühe silmaga läbi pliiatsite vahelise pilu ja kirjeldage, mida on võimalik näha. Püüdke nähtut seletada.

\hint

\solu
Pilu keskel on tume kriips --- 1 p. Küünla valgus venib kahele poole välja --- 1 p. Tekivad tumedad kohad väljavenitatud osasse --- 1 p. Tumedate osade vahelised alad on värvilised --- 1 p. Keskkohale lähemal olevad servad on sinised, kaugemad punased --- 1p. Kui muuta pilu kitsamaks (kas pööramise või pigistamisega), siis pilt venib laiemaks --- 2 p. Pilu pööramisel ümber vaatesuuna pöördub ka valgusriba --- 1 p. Pilt oleneb silma kaugusest pilust --- 1 p. Pilt oleneb pilu kaugusest küünlast ja on paremini nähtav suuremate kauguste korral --- 2 p. Nähtust nimetatakse difraktsiooniks, mis seisneb valguse kandumises varju piirkonda --- 1 p. Pilust väljuvad valguslained võivad liitudes teineteist tugevdada või nõrgendada --- 1 p.
\probend
